\documentclass{ximera}

\title{Logarithm Review Activity}
\author{MATH 425: Calculus I}

\begin{document}
\begin{abstract}
   Working with the peers in your group, solve the following problems. Make sure to show and justify all your work. Make sure everyone in the group understands the solution and participates. Be prepared to report your answers to the whole class. 
\end{abstract}
\maketitle

\begin{exercise}
    Rewrite the following logarithmic expressions as exponential functions, and exponential functions as logarithmic functions.  In each case, your answer should be an equation with a single logarithm or a single exponential function on one side of the equation.
    \begin{enumerate}
        \item $\log_5 25=2$
        \item $\log_2 8=3$
        \item $\log_{10} 1000=3$
        \item $e^3\approx 20.0855$
        \item $10^4=10000$
    \end{enumerate}
\end{exercise}


\begin{exercise}
    For a population that is growing exponentially according to a model of the form $P(t)=Ae^{kt}$, the doubling time is the amount of time that it takes the population to double.  For each population described below, assume the function is growing exponentially according to the model $P(t)=Ae^{kt}$, where $t$ is measured in years.  
    \begin{enumerate}
        \item Suppose a certain population initially has 100 members and doubles after 3 years.  What are the values of $A$ and $k$ in the model?
        \item A different population is observed to satisfy $P(4)=250$ and $P(11)=500$.  What is the population's doubling time?  When will 2000 members of the population be present?
        \item Another population is observed to have doubling time $t=21$.  What is the value of $k$ in the model?
        \item How is $k$ related to the doubling time, regardless of how long the doubling time is?
    \end{enumerate}
\end{exercise}

\begin{exercise}
    \begin{enumerate}
        \item Find the solution of the exponential equation $1000(1.04)^{2t}=50000$ in terms of logarithms. (Keep your answer exact.)
        \item Find the solution to the logarithmic equation $2-\ln(3-x)=0$ correct to three decimal places.
    \end{enumerate}
    
\end{exercise}

\begin{exercise}
    Our goal in this exercise is to investigate the product property for logarithms.
    \begin{enumerate}
        \item Write $10^x\cdot 10^y$ as $10$ raised to a single power.  That is, complete the equation \[10^x\cdot 10^y = 10^{\square}\]
        by filling in the box with an appropriate expression involving $x$ and $y$.
        What is the simplest way to write $\log_{10}10^x$?  What about the simplest equivalent expression for $\log_{10}10^y$?
        Explain why each of the follwoing three equal signs is valid in the sequence of inequalities:
        \begin{align}
            \log_{10}(10^x\cdot 10^y) & =\log_{10}(10^{x+y})\\
            &=x+y\\
            &= \log_{10}(10^x) + \log_{10}(10^y)\\
        \end{align}
        \item Suppose that $a$ and $b$ are positive real numbers so we can think of $a$ as $10^x$ for some real number $x$ and $b$ as $10^y$ for some real number $y$. That is, $a=10^x$ and $b=10^y$. What does our work in (c) tell us about $\log_{10}(ab)$?
    \end{enumerate}
\end{exercise}

Problems 2-4 from Boelkins, et al., \textit{Active Prelude to Calculus}.

\end{document}
